\section{Discussion and Limitations}
\label{sec:limitations}

\paragraph{Backbone depth.} Every limitation observed empirically traces back to Conv4's shallow, low-resolution final feature grid: it widens \biadt's confidence margins (\Cref{tab:biadt}), causes spatial bleeding onto background for small objects (\Cref{fig:case_study}), and leaves the cascading-randomization sweep (\Cref{tab:sanity}a) too coarse (only 5 destruction levels) to resolve exactly \emph{where} within the collapsing block the useful gradient signal disappears. None of these change the sign or significance of any result; they are resolution artifacts, not faithfulness failures, and we expect deeper backbones (e.g.\ ResNet-12) to sharpen all three.

\paragraph{Single configuration, no explainer baseline.} Every result in \Cref{sec:experiments} uses one algorithm/backbone pair, MAML/Conv4; we have not yet re-run \Cref{tab:biadt,tab:sanity} on a second backbone or algorithm, so cross-architecture generality is argued qualitatively (\Cref{sec:related}), not shown quantitatively. More importantly, we validate that $\PhiFama$ is faithful and sanity-passing in an absolute sense, but do not yet compare it head-to-head, under the same \biadt protocol and on the same tasks, against the closest competing attributions: static Grad-CAM applied only at $\phi_T$ (isolating how much of \fama's faithfulness comes specifically from aggregating the full trajectory, \Cref{eq:branch_saliency}, rather than from the sign-preserving CAM projection alone), TL-XML~\cite{mitsuka2025tlxml} (task-level rather than pixel-level, and requiring meta-training data we do not need), and a TracIn-style~\cite{pruthi2020estimating} trajectory sum without the adjoint's Hessian curvature term. Without this comparison, \Cref{sec:experiments} shows \fama is internally faithful, not that it is more faithful than a cheaper static alternative -- the strongest open question against the method as it stands. None of the three baselines requires new machinery: all operate on the same trajectories and checkpoints \fama already consumes, so this -- together with a wall-clock comparison against the pass counts in \Cref{sec:fama} -- is a well-scoped follow-up rather than an open problem, gated on compute rather than on anything unresolved in the method itself.

\paragraph{OOD sensitivity.} The OOD support-set intervention (\Cref{tab:sanity}b) is the weakest of our three sanity signals -- not because $\PhiFama$ is insensitive to domain shift, but because CUB-200-2011 and tieredImageNet, despite being semantically disjoint, share enough low-level visual statistics (centered subject, contrasting background) to still activate Conv4's shallow filters similarly. A non-object OOD source (e.g.\ textures or scenes) should make this a sharper stress test.

\paragraph{Architectural scope.} \fama's upsampling step (\Cref{eq:branch_saliency}) assumes the feature grid preserves spatial correspondence with the input, as in standard CNNs; extending it to architectures that do not preserve this correspondence (e.g.\ pure Transformers) is not automatic and is left to future work, together with validating $\Delta M$ and \fama across other members of the MAML family (ANIL, BOIL, Meta-SGD).

\paragraph{From diagnosis to curation.} \Cref{sec:fama} already produces what a concrete data-curation rule would need: a signed, per-pixel map $\PhiFama(x_k^S)$ for every support image. A direct proposal is to reduce this to a per-image score $s_k := \sum_{i,j}\PhiFama(x_k^S)_{i,j}$, then either \emph{down-weight} $x_k^S$'s gradient contribution in \Cref{eq:inner_loop} in proportion to $s_k$ (rather than discarding it outright, so $S_i$ never loses class coverage), or \emph{flag} the lowest-$s_k$ images in $S_i$ for human review before deployment. Both reuse the maps \fama already computes; whether either one measurably recovers query performance on the negative-transfer tail of \Cref{fig:delta_m_dist}, and how it compares to coupling \fama with a post-hoc representation editor such as HSFM~\cite{parast2026hsfm}, is future work.
